\section{Origin of the Race of the Soul}

In formulating a theory of the origin of the race of the soul, Julius Evola points out that ``soul'' (or psyche) does not mean ``subjective activities'' in the mind as contemporary psychology understands it. Rather, it is a distinct entity, which he relates to the traditional teachings of the ``subtle body''. In other words, it is not reducible to material, physical, or biological factors. Hence, scientific methods cannot be applied to understand the race of the soul.

At its origins, the soul and the body may have been in harmony, but that is seldom the case today. Hence, there are really two levels of heredity: the biological and the psychical. Moreover, psychical heredity is non-deterministic since the factor of elective affinity comes into play.

\paragraph{Addendum}


A commenter mentioned that this chapter seemed ``arbitrary and even dangerous.'' However, I believe this chapter is significant, at least in some aspects, if not in toto. For example, I believe it offers an explanation for the transmission of what is called ``original sin''. In times past, the method of transmission was debated, e.g., whether it came through semen, etc. However, such biological and material explanations make no sense to anyone today. Were that true, there would be an ``original sin'' gene. Then, presumably, gene therapy could someday reverse the effects of original sin.

On the other hand, Joseph Ratzinger's description of original sin as more like a psychological matrix that we are unavoidably born into makes more sense to people today. It can be understood and even observed. Nevertheless, it is still weak since a means of transmission is never specified. The notion that there is a psychical heredity, therefore, answers the question of transmission. Moreover, psychical deformations even spill over onto the body, with ill effects. This is just a starting point for understanding that doctrine.

\paragraph{Postscript}


The German translation has some interesting additions. In it, Evola refers to the ``Nordic-Roman'' race rather than the ``Nordic'' race.

He also quotes Alfred Rosenberg, with whom he otherwise had significant disagreements:

\begin{quotex}
The body is the externality of the soul and the soul is the race seen from the inside.
\end{quotex}

\paragraph{Origin of the Races of the Soul}


Where do the races of the soul come from? Clearly, in the limit-case of completely pure races from a single stream, if one could assume it, they represent the psychological expression of the same particular formative power that is expressed on the physical plane by the specific and typical characteristics of the anthropological race of the body which forms the base of their inseparable unity, while belonging in itself to a still higher plane. According to the ancient traditional teachings, the soul is not simply what modern psychology believes, that is to say, an ensemble of phenomena and ``subjective'' activities, unfolding on a physico-biological base. For that teaching, the soul is rather a type of being in itself; as the \emph{linga sharira}, or ``subtle body'', it has its own existence, its real forces, its laws, its own heredity, distinct from the purely physico-biological.

From such a point of view, it is necessary to consider that the races of the soul are subject to vicissitudes similar the race of the body, except that, to indicate such vicissitudes, and therefore in order to know about the genesis of the races of the soul, about their essence and laws as conditioning their development and their integrity, it would necessarily be by means of immaterial research previously known by the ancient traditional sciences, but now unknown to modern culture. Although a deformed recollection of it is found in certain theosophical and occultist circles, there is however not even a suspicion of it in so-called ``scientific'' research. So as things stand today, it is necessary to proceed in an inductive or intuitive way, rather than from a clear body of knowledge.

In any case, bear in mind the important methodological foundation of the principle that two distinct lines of heredity exist, one of the body and the other of the soul, lines which may even diverge after races and traditions have lost their original purity from prehistoric times. While the line of biological heredity has a visible and identifiable continuity because it is based on the process of natural generation, we have to assume that the line of the heredity of the soul instead has its continuity only on another, no longer visible, plane and can therefore connect individuals who might have nothing in common in space and time. We will return to this point in discussing the problem of birth. Here we observe the complexity that previously the problem of physical heredity itself presents in these terms, even if it not is considered as positivistic myopia. In fact, since the soul has a relationship of reciprocal action with the body, in the case of the divergence of the two heredities, there will arise in physical heredity, through the influence of the other, modifications insusceptible to explanations that biological and anthropological research will never be able to ascertain within its domain.

However, this is not the suitable place to discuss such topics, because it presupposes knowledge of the traditional conception of the multiple states of being, that substitute for the way that all the greatest problems related to man, life, death, and the world are considered today. Let's return to the point from which we started in order to say that, when you are faced with a state of racial mixing, the races of the soul should be considered as the result of three factors. The first, which is essential for it, is precisely the race of the soul as a distinct entity. The second is the influence that a mismatched body of race may have exercised over it and, through this body, which is the positive center of the relations with the outer world, an unsuitable environment. The third is the influence that a still higher element, i.e., the race of the spirit, may have exercised in the case of a new divergence between the soul and spirit, beyond that between the soul and body.

Since the unity of the various elements does not happen by chance and automatic laws, but on the basis of analogous and ``elective'' connections (this will also be clarified later), in spite of the divergences, we can admit, as a working hypothesis and a likely criterion, a certain equivalence. For example, out of the hundred types that demonstrate racial purity as race of the body, let us suppose, of the Nordic type, we can presume that a greater number of cases are found where, a psycho-spiritual congruent qualification also virtually corresponds to it, than between a hundred types whose race of the body is neither Nordic, nor of Nordic origin. The misgivings of making such assumptions follow. First of all, the one just made by saying ``virtually'', since, as we saw, there were cases of pure races that are half-extinguished or have declined in terms of the race of the soul. In the second place, because it is necessary to consider the cases of ``propensity'' -- the law of affinity may have made the manifestation of a given type of personality prefer a certain race of the soul, however in accordance with such a circumstance, made to pay for this choice with the reception of a non-matching race of the body (for example, in cases of reanimation of the race in the second of the forms considered in part 2, chapter 11, the elective affinities would even lead to a manifestation in mixed forms more so than in pure, but interiorly decadent, forms). In the third place, because ``analogy'' and ``elective affinity'' are terms that refer to not simply human states of existence, so that through it criteria apply which may not match the way the common mind has been led to suppose and to believe natural, logical, and desirable.


\flrightit{Posted on 2015-05-17 by Aeneas}

\begin{center}* * *\end{center}

\begin{footnotesize}
\begin{sffamily}
\texttt{Paulo on 2015-05-17 at 23:08 said:}

Other important aspect that should be demonstrated,is the influence of the land that we are born.In Carl Jung's book,''civilization in transition''in chapter XXII he mention a certain professor Boas,that have made studies that shows that the sons of european imigrants born in the United States,had a different anatomy then their ancestors.Later in the same chapter,Jung mention some primitive people that belives that the children born in the land usurped by their fathers from native americans,africans or whatever be the case,shoul''inherit the spirits''of the land.He see some truth in that!!!

In the same chapter Jung writes extensively about the psychological differences between americans and their europeans counterpart!!!

\hfill

\texttt{Jacob on 2015-05-18 at 13:26 said:}

Thank you for the translations.

How does this relate to Tomberg's reincarnation as psychological fact or does it?

\hfill

\texttt{Cologero on 2015-05-18 at 16:35 said:}

That's a good point, Jacob. The two Hermetists agree that there is biological and psychic heredity (bracketing out spiritual heredity for the moment). Tomberg discusses this in Letter X. However, Tomberg relates psychic continuity to ``reincarnation'', while Evola denies reincarnation. It probably comes down to a verbal disagreement. Tomberg denies the crude teaching of reincarnation in which ``I'' am being punished for something ``I'' did in the past.

What he is asking is that if ``I'' am Peter and ``I'' have a memory of Paul, does that mean I am both Peter and Paul? It can't mean that Paul reincarnated as Peter, since the ``I'' is transcendent to the experience of either one.

Nevertheless, Tomberg wants to call that an ``identity'' of some sort. Evola expresses it differently: Peter has an ``affinity'' with Paul, so Peter is born within the psychic current that included Paul. In other words, the experience may be the same, but the concept of affinity explains it better than the concept of identity. Guenon objects that an ``identity'' means that two things are the same in every aspect, but clearly Peter and Paul are different.

Coomaraswamy's essay on the One and Only Transmigrant may be helpful. Ultimately, there is the One I, manifesting in different ways.

\end{sffamily}
\end{footnotesize}

